\project{07}{Stop mode, RTC wake, and a battery number}{NUCLEO-H7A3ZI-Q with an nRF-PPK2}{Low-power modes, clock restore after wake, charge per cycle}

\begin{keyfacts}
\item[Board] NUCLEO-H7A3ZI-Q, with an nRF-PPK2 in the supply path
\item[Peripherals] PWR and RCC for the mode and the clock restore, the RTC on the fitted 32.768 kHz crystal, LPTIM1, EXTI for the wake line, three GPIO markers for the instrument's digital channels
\item[Toolchain] arm-none-eabi-gcc with CMake as in chapter 1; Python 3 on the host with the instrument's interface installed from its package index
\item[Operating system] Bare metal
\item[Difficulty] 5 of 5
\item[Effort] 4 evenings of about four hours, and the last one is entirely arithmetic
\item[Deliverable] A duty-cycled application that sleeps between wakes, restores its clock tree correctly, and reports charge and energy per cycle with a spread over at least ten cycles, together with an honest statement of what this part is and is not good for
\end{keyfacts}

\subsection*{Why this project}

This chapter has to begin with the finding, because the finding changes what
the chapter is allowed to claim. \textbf{This part is a poor low-power node.}
The literature on energy-budgeted sensor nodes runs on the ultra-low-power
families: the STM32L4 and STM32U5 lines, the nRF52 and nRF54 series, the
MSP430. Those parts idle in microamps. This one idles in milliamps, on a board
whose debug probe is powered from the same cable. A search of the published
work turns up no thesis and no paper that builds a low-power sensor node on this
family at all, which is a negative result worth recording rather than a gap to
apologise for.

So the chapter keeps its subject and changes its claim. It is not "here is a
low-power node". It is "here is how you account for energy, and here is what
racing to sleep looks like on a part chosen for speed". Both halves are
transferable. The accounting method is the same on any part, and it is the half
that most published low-power writing does badly: an average current with no
spread, taken over an unstated interval, on a board with unstated peripherals
enabled. The racing-to-sleep half is specific and interesting here precisely
because the part is fast: a processor that finishes its work in a tenth of the
time can afford a much worse sleep current before it loses, and the crossover is
a number this chapter can produce.

There are two pieces of good news, both verified rather than assumed. The
32.768 kHz crystal is fitted on this board. That is confirmed from two
independent machine-readable sources that name this exact board: the upstream
board description in the other operating system's tree lists a 32.768 kHz
crystal oscillator, and the MicroPython board definition selects that oscillator
as the real-time clock source. So the real-time clock is usable here and its
accuracy is claimable, which is what makes a ten second wake period a period
rather than an approximation. And the vendor's firmware package carries three
power examples for this exact board under
\texttt{Projects/}\allowbreak\texttt{NUCLEO-H7A3ZI-Q/}\allowbreak\texttt{Examples/PWR}: standby, standby with a clock
wake, and the deeper stop mode with a clock wake. Working register sequences
for this board exist and do not have to be derived from a sibling part.

\begin{note}{The clock tree does not come back on its own}
This is the trap specific to this family and it is the reason the chapter has a
whole step for it. On wake from the deeper sleep state the system clock has
reverted to the internal oscillator and the phase-locked loop is off. Nothing
restores it. The symptom is not a failure: the application runs, the wake
completes, and the measured current after the wake stays higher than it should
while the core runs slower than it should. Every project that reports
disappointing energy numbers on this family should check this first. The
restore also has an order: the voltage scale and the flash wait states come
before the frequency, exactly as in chapter 1.
\end{note}

\subsection*{Prior art and what to reuse}

\begin{table}[H]\centering\small
\begin{tabular}{p{32mm}p{46mm}p{40mm}p{20mm}}
\toprule
Source & What it gives & What it does not & Licence \\
\midrule
The vendor's three power examples for this exact board & Working entry and wake sequences for standby, standby with a clock wake, and the deeper stop mode with a clock wake, written against RM0455 & No measurement, no clock restore worth the name, no accounting, and nothing about what a cycle costs & Basic examples are BSD-3-Clause; read the header of each file, because the middleware around them is not \\
The instrument's Python interface & Scripted control of the power instrument: mode, supply voltage, start and stop, and the sample stream with its digital channels & It is a transport, not a method. It computes no charge and segments no cycles & GPL-2.0. Install from the package index. Do not fork it into a portfolio repository \\
Golioth's article on measuring current with this instrument & The procedure: how to place the instrument, which mode to choose, and how to read a trace & It is an introduction on a different board, with no uncertainty treatment and no per-cycle statistics & Article, read and cite \\
The MicroPython board definition for this board & Confirms the low-speed crystal is fitted and is used as the real-time clock source & It is an interpreter, and its power behaviour is not this chapter's subject & MIT \\
The ultra-low-power families' datasheets and application notes & The comparison to beat, in numbers their vendors publish and stand behind & They are different silicon. Their figures are the target, never this part's result & Vendor documentation \\
\bottomrule
\end{tabular}
\caption{Prior art for chapter 7. No peer-reviewed method exists for this instrument, so the measurement rigour below is original work rather than a reproduction.}
\end{table}

What is left to write is the rigour and the restore. The clock restore
function, the phase marking that lets the trace be segmented automatically, the
per-cycle charge computation with a spread over at least ten cycles, and the
separation of the measurement from the battery arithmetic are the author's, and
chapters 10 and 12 use all four unchanged.

\subsection*{Parts from the inventory}

\begin{table}[H]\centering\small
\begin{tabular}{p{40mm}p{66mm}p{40mm}}
\toprule
Part & Role & Interface \\
\midrule
NUCLEO-H7A3ZI-Q & The subject. Duty cycles between a sense phase and a sleep state & Micro USB to the host \\
nRF-PPK2 & Ammeter or source meter up to 1 A, with eight digital inputs time-aligned to the current trace & USB to the host, two leads to the board \\
Three jumper wires & Two for the supply path, one common ground for the digital channels & Jumper header \\
Three more jumper wires & Phase markers from free GPIO to the instrument's digital inputs & Zio header to the instrument \\
Host PC & Runs the capture script and the analysis & USB \\
\bottomrule
\end{tabular}
\caption{Inventory items used in chapter 7. The instrument measures at most 1 A, which is comfortable for this board but is the reason a cellular modem is never placed in the same path; chapter 11 says why.}
\end{table}

\subsection*{System architecture}

\diagram{c07_arch}{The instrument sits in the supply path and watches three marker pins at the same time. Everything the chapter claims comes out of one trace with four channels in it.}

The instrument is in series with whatever it is measuring, which is the first
thing that has to be settled and is not settled yet. Two topologies are
possible on this board. In ammeter mode the instrument is placed across the
jumper that carries the microcontroller's own supply current, and the board is
powered as usual. In source-meter mode the instrument supplies the board
itself. Which jumper carries that current, exactly what it isolates, whether
the probe's own current is excluded, and whether the instrument can power the
board through it at the required voltage are items 6 and 34 on the confirm
list. Until they are settled the chapter measures in ammeter mode and states in
every caption that the probe is still powered from USB and is not in the
measured path.

\subsection*{Peripheral configuration}

\begin{table}[H]\centering\small
\begin{tabular}{p{20mm}p{28mm}p{24mm}p{30mm}p{30mm}}
\toprule
Peripheral & Mode & Clock source & Pins and function & Interrupt and transfers \\
\midrule
PWR & Stop with the deepest stop voltage scale the part allows & None & None & None \\
RCC & Restore path after wake: voltage scale, flash latency, phase-locked loop, switch & 8 MHz bypass from the probe & None & None \\
RTC & Calendar and a periodic alarm & The fitted 32.768 kHz crystal & None & Alarm through the wake line \\
LPTIM1 & Periodic counter that continues in the stop state & The same crystal & None & Update through the wake line \\
EXTI & Wake line for the alarm and for the user button & None & PC13 for the button & Wake from stop \\
GPIO markers & Push-pull outputs, one per phase & Peripheral bus & Three free Zio pins, item 13 on the confirm list & None \\
SysTick & Stopped before sleeping, restarted after & Core clock & None & Tickless build only \\
\bottomrule
\end{tabular}
\caption{Peripheral configuration. What each low-power mode retains, which sources can wake from it and what the wake latency is are item 27 on the confirm list and are read out of RM0455's low-power section before any of them is printed as fact.}
\end{table}

\subsection*{Wiring}

\diagram{c07_wiring}{The instrument in ammeter mode, in series with the board's own supply through the measurement jumper, with the three marker pins on its digital inputs and one common ground. Which jumper this is remains item 6 on the confirm list.}

Three safety points. The instrument measures at most 1 A, so nothing that draws
a pulse larger than that goes in this path, which excludes every cellular modem
on this bench. The board's logic is 3.3 V and the instrument's digital inputs
are driven by the board, never the other way round; the logic threshold of
those inputs is item 33 on the confirm list and is read before the markers are
trusted. And the ground between the instrument and the board must be common or
the digital channels float, which produces a trace that looks plausible and
means nothing.

\subsection*{Memory and timing budget}

\diagram{c07_mem}{The clock tree at the instant of wake, next to the clock tree the application expects. The difference between them is the whole of the trap, and nothing in the silicon closes it.}

\begin{table}[H]\centering\small
\begin{tabular}{p{44mm}p{28mm}p{28mm}p{30mm}}
\toprule
Quantity & Budget & Measured & Margin \\
\midrule
Wake period & 10.000000 s from the fitted crystal & not measured & not measured \\
Run current at 280 MHz & no budget until the first measurement & not measured & not measured \\
Stop current & no budget until the first measurement & not measured & not measured \\
Wake to first instruction & under 1 ms & not measured & not measured \\
Wake to restored clock tree & under 2 ms & not measured & not measured \\
Charge per cycle & established here & not measured & not measured \\
Energy per cycle & established here & not measured & not measured \\
Spread over ten cycles & within 5 percent of the median & not measured & not measured \\
\bottomrule
\end{tabular}
\caption{The budget for chapter 7. The current rows carry no budget on purpose. A budget invented before the first measurement becomes the number the measurement is read against, and no figure in microamps is written anywhere in this book until an instrument has produced it.}
\end{table}

One row in that table is exact and it is worth saying why. The wake period is
ten seconds because the crystal is 32.768 kHz and the real-time clock's
prescalers divide it by exactly 32768, so a ten second alarm is ten seconds and
not ten seconds plus a drift term. That is only true because the crystal is
fitted, which is a confirmed fact about this board and not a family assumption.
On a board where the low-speed oscillator is the internal one, the same code
produces a period with a tolerance of several percent and every per-cycle figure
inherits it.

\subsection*{Firmware design (UML)}

\diagram{c07_uml}{The duty cycle as a state machine, with the state most projects on this family omit drawn in full. The path that skips the restore is drawn too, because it is the one that runs.}

The application is a loop with a sleep in it, and the only design decision that
matters is where the boundaries of a cycle are. This chapter puts them at the
wake, so a cycle is wake, restore, sense, prepare, sleep, and the next wake. The
marker pin that encloses the whole cycle is what lets the analysis segment the
trace without a human looking at it, and that is the difference between a
measurement you can repeat ten times and a screenshot.

\subsection*{Data flow (ASCII)}

\begin{asciiart}
  board                                            instrument                host
  +-------------------------------------+          +---------------+         +------------------+
  | RTC alarm  (32.768 kHz crystal)     |          | nRF-PPK2      |         | capture.py       |
  |   |                                 |          |               |         |   mode, voltage  |
  |   v                                 |   IDD    | ammeter in    |  USB    |   start, stop    |
  | wake: HSI, PLL off, slow            +--------->| series with   +-------->|                  |
  |   |                                 |          | the supply    |         | cycles.py        |
  |   v                                 |          |               |         |   segment by D0  |
  | restore: VOS, latency, PLL, switch  |          | D0 cycle      |<--------+   integrate I dt |
  |   |                                 |  D0 D1 D2| D1 sense      |         |   per-cycle Q    |
  |   v                                 +--------->| D2 restore    |         |   median, spread |
  | sense: one sample, one feature      |          |               |         |                  |
  |   |                                 |          | 4 channels,   |         | battery.py       |
  |   v                                 |          | one timebase  |         |   Q -> lifetime  |
  | prepare: markers low, tick stopped  |          +---------------+         |   assumptions    |
  |   |                                 |                                    +------------------+
  |   v                                 |
  | stop: SLEEPDEEP, WFI                |
  +-------------------------------------+
\end{asciiart}

\subsection*{Repository layout}

\begin{plaincode}
nucleo-h7a3-lowpower/
  CMakeLists.txt                  # three sleep strategies from one application
  cmake/arm-none-eabi.cmake       # unchanged from chapter 1
  src/main.c                      # the duty cycle, identical in all three
  src/sleep.h                     # the one interface all three satisfy
  src/sleep_rtc.c                 # wake on a real-time clock alarm
  src/sleep_lptim.c               # wake on the low-power timer
  src/sleep_tickless.c            # stop the tick, sleep to the next deadline
  src/clocks.c                    # restore_after_stop(), the chapter's core
  src/markers.c                   # three phase pins, shared with chapter 10
  meas/capture.py                 # instrument control, writes a raw trace
  meas/cycles.py                  # segment, integrate, per-cycle statistics
  meas/battery.py                 # the arithmetic, and what it assumes
  meas/test_cycles.py             # the analysis tested on a synthetic trace
  meas/requirements.txt           # the interface pinned, never vendored
  docs/method.md                  # written before the first capture
  README.md
\end{plaincode}

\subsection*{Steps}

\begin{steps}
\item \textbf{Install the instrument's interface without vendoring it.} The
Python interface to this instrument is GPL-2.0. That licence is entirely fine
to depend on and entirely wrong to copy into a portfolio repository that you
intend to publish under a permissive licence. Install it, pin it, and record
why in the README.
\begin{shellcode}
python -m venv .venv && . .venv/bin/activate
pip install ppk2-api
pip freeze > meas/requirements.txt
\end{shellcode}

\item \textbf{Prove the instrument before you believe it.} An instrument that
has not been checked against something known is an opinion. Put a resistor of
known value across the supply at a known voltage and confirm the reading is
where Ohm's law puts it, then record the deviation in the README as the
instrument's working uncertainty for the rest of this book.
\begin{pycode}
# meas/capture.py -- the API names follow the published examples of the
# version pinned in requirements.txt. Check them against that version.
from ppk2_api.ppk2_api import PPK2_API

ppk = PPK2_API("/dev/ttyACM0")
ppk.get_modifiers()                 # calibration constants from the device
ppk.use_ampere_meter()              # the board keeps its own supply
ppk.start_measuring()
raw = ppk.get_data()
current_uA, digital = ppk.get_samples(raw)
ppk.stop_measuring()
\end{pycode}
The digital channels come back alongside the current samples on one timebase,
which is the property this whole chapter is built on. How many channels, at what
logic threshold, at what sample rate, and whether an exported file carries them
are item 33 on the confirm list.

\item \textbf{Measure the run state first and write down what was enabled.} A
current figure with no statement of which peripherals were on is not a
measurement. Record the core frequency, the voltage scale, which peripheral
clocks are enabled, whether the caches are on, and whether the debugger is
attached. Attach that list to the number, permanently.

\item \textbf{Enter the stop state with a real-time clock alarm.} The domain
structure of this part differs from its better-known sibling: it has a CPU
domain and a smart-run domain where the other has three power domains, so the
bit names in the power control registers are not the same. Take the sequence
from RM0455's low-power section and from the vendor's example for this exact
board.
\begin{ccode}
/* sleep_rtc.c: arm the alarm, then stop. Register names from the CMSIS
   device header for this part; the deepest stop voltage scale the part
   allows is read out of RM0455, not assumed from a sibling. */
void sleep_until_next_alarm(uint32_t seconds)
{
    rtc_set_alarm_in(seconds);              /* alarm A, through the wake line */

    PWR->CR1 = (PWR->CR1 & ~PWR_CR1_SVOS_Msk) | PWR_SVOS_DEEPEST;
    SCB->SCR |= SCB_SCR_SLEEPDEEP_Msk;      /* stop, not plain sleep */

    marker_clear(MARK_CYCLE);               /* the trace's cycle boundary */
    __DSB();
    __WFI();                                /* execution resumes on the line
                                               after this one, at reduced
                                               speed, with PLL1 off */
    marker_set(MARK_CYCLE);
    SCB->SCR &= ~SCB_SCR_SLEEPDEEP_Msk;
}
\end{ccode}
Two details decide whether this works at all. Every wake source must be enabled
in the wake-line register as well as in its own peripheral, because the
peripheral interrupt alone does not leave the stop state. And any pending
interrupt at the moment of the wait instruction causes it to return
immediately, which presents as a board that never sleeps and is almost always
a flag left set from the previous cycle.

\item \textbf{Restore the clock tree, and measure the difference.} This is the
step the chapter exists for. Run one capture with the restore commented out and
one with it in, and put the two traces side by side.
\begin{ccode}
/* clocks.c: the order is voltage scale, then wait states, then frequency.
   Reversing it is the classic way to produce a board that runs at reset
   speed, or one that does not run. */
void restore_after_stop(void)
{
    marker_set(MARK_RESTORE);

    pwr_set_run_voltage_scale();            /* back to the scale 280 MHz needs */
    flash_set_latency_for_280mhz();         /* wait states BEFORE the frequency */

    RCC->CR |= RCC_CR_HSEON;                /* bypass input from the probe */
    while ((RCC->CR & RCC_CR_HSERDY) == 0u) { }
    RCC->CR |= RCC_CR_PLL1ON;
    while ((RCC->CR & RCC_CR_PLL1RDY) == 0u) { }

    MODIFY_REG(RCC->CFGR, RCC_CFGR_SW, RCC_CFGR_SW_PLL1);
    while ((RCC->CFGR & RCC_CFGR_SWS) != RCC_CFGR_SWS_PLL1) { }

    marker_clear(MARK_RESTORE);
}
\end{ccode}
The restore costs time and the marker measures it. That cost is what makes the
racing-to-sleep argument quantitative rather than rhetorical: a part that takes
longer to become fast again has a longer fixed overhead per cycle, and below
some wake period the overhead is the whole budget.

\item \textbf{Build the low-power timer variant.} The general timer of chapter 6
stops in the stop state because its bus clock stops. The low-power timer does
not, because it runs from the same crystal the real-time clock uses and its
logic continues. That is the entire reason it exists and it is why this chapter
rather than chapter 6 builds it.
\begin{ccode}
/* sleep_lptim.c: a periodic wake without the calendar. */
void lptim_start_periodic(uint32_t ticks)
{
    /* kernel clock select: the fitted 32.768 kHz crystal, not the internal
       low-speed oscillator. The selection field is in RCC's kernel clock
       configuration register for this part. */
    rcc_select_lptim1_kernel_clock_lse();

    LPTIM1->CR  = 0u;
    LPTIM1->IER = LPTIM_IER_ARRMIE;         /* autoreload match */
    LPTIM1->CR  = LPTIM_CR_ENABLE;
    LPTIM1->ARR = ticks - 1u;               /* write only while enabled */
    LPTIM1->CR |= LPTIM_CR_CNTSTRT;         /* continuous */
}
\end{ccode}
The autoreload register on this peripheral is written after the peripheral is
enabled and the write takes time to be accepted. Writing it before enabling, or
writing it again without waiting, is the standard first fault with this
peripheral and it presents as a period that is right sometimes.

\item \textbf{Build the tickless variant.} The millisecond tick is a source of
wakes that do nothing. In this build the tick is stopped before sleeping and
the elapsed time is reconstructed on wake from the counter that kept running.
\begin{ccode}
/* sleep_tickless.c: stop the tick, sleep, then account for the gap. */
void sleep_tickless(uint32_t ms)
{
    uint32_t before = lptim_count();
    SysTick->CTRL &= ~SysTick_CTRL_ENABLE_Msk;

    lptim_set_next_match(before + ms_to_ticks(ms));
    sleep_until_next_alarm_or_match();

    uint32_t slept = lptim_count() - before;        /* unsigned, wrap is free */
    tick_advance(ticks_to_ms(slept));               /* the clock did not stop */
    SysTick->CTRL |= SysTick_CTRL_ENABLE_Msk;
}
\end{ccode}
Two things are worth stating plainly. The reconstruction has a resolution of one
crystal tick, which is about 30.5 \textmu{}s, so a millisecond software clock
built this way is accurate but granular. And the same scheme on an operating
system is that system's tickless idle, which chapter 20 builds; doing it by hand
first is what makes that chapter's version legible rather than magical.

\item \textbf{Mark the phases and segment the trace automatically.} Three
markers: one enclosing the whole cycle, one on the sense phase, one on the
restore. The analysis never asks a human where a cycle begins.
\begin{pycode}
# meas/cycles.py
import numpy as np

def per_cycle_charge(current_uA, digital, fs, cycle_bit=0):
    """Charge in microcoulombs for each complete cycle in the trace."""
    d = (digital >> cycle_bit) & 1
    edges = np.flatnonzero(np.diff(d) != 0)
    rises = edges[d[edges + 1] == 1]
    out = []
    for a, b in zip(rises[:-1], rises[1:]):
        out.append(current_uA[a:b].sum() / fs)      # uA * s = uC
    return np.asarray(out)

def report(q_uC, volts=3.3):
    e_uJ = q_uC * volts
    return {
        "cycles":        int(q_uC.size),
        "charge_median": float(np.median(q_uC)),
        "charge_min":    float(q_uC.min()),
        "charge_max":    float(q_uC.max()),
        "energy_median": float(np.median(e_uJ)),
        "spread_pct":    float(100 * (q_uC.max() - q_uC.min())
                               / np.median(q_uC)),
    }
\end{pycode}
Report the median and the extremes over at least ten cycles, never a single
average current. An average current over a duty-cycled trace hides the thing
that decides a battery life, which is whether the cycles are alike.

\item \textbf{Test the analysis on a synthetic trace.} The segmentation is the
fragile part, so it is tested on the host with a trace whose answer is known by
construction, in continuous integration, with no instrument attached.
\begin{pycode}
# meas/test_cycles.py
import numpy as np
from cycles import per_cycle_charge

def test_known_charge():
    fs = 100000.0
    n_cycle = int(fs)                       # one second per cycle
    i = np.full(10 * n_cycle, 100.0)        # 100 uA flat, 10 cycles
    d = np.tile(np.r_[np.ones(n_cycle // 10), np.zeros(9 * n_cycle // 10)], 10)
    q = per_cycle_charge(i, d.astype(np.uint8), fs)
    assert q.size == 9                      # nine complete cycles in ten
    assert np.allclose(q, 100.0, rtol=1e-6) # 100 uA for 1 s is 100 uC
\end{pycode}

\item \textbf{Turn the measurement into a battery number and keep the two
apart.} The energy per cycle is a measurement. The lifetime is arithmetic over
that measurement and a set of assumptions, and the assumptions are where these
claims usually fail. State them next to the number, every time.
\begin{pycode}
# meas/battery.py
def lifetime_days(energy_per_cycle_uJ, cycles_per_day, cell_mAh, cell_V):
    """Arithmetic only. The measurement is the first argument; everything
    else is an assumption and is printed with the answer."""
    cell_J = cell_mAh * 1e-3 * 3600.0 * cell_V
    return cell_J / (energy_per_cycle_uJ * 1e-6 * cycles_per_day)
\end{pycode}
What that arithmetic does not model, and what the README says it does not model:
self-discharge, the capacity a coin cell actually delivers under pulsed load
rather than under its rated continuous load, temperature, the regulator's
efficiency at the currents involved, and the fact that this board is not the
product. Whether the backup domain is usable without a coin cell at all is item
10 on the confirm list. A battery number offered without that paragraph is
advertising rather than engineering.
\end{steps}

\subsection*{Build, flash and debug}

\diagram{c07_timing}{One cycle with the restore and one without, on the same axes, with the marker channel underneath. The shape of the difference is the point; the values on the vertical axis come from the instrument and are not drawn here because they have not been measured.}

\begin{shellcode}
cmake -B build -DCMAKE_TOOLCHAIN_FILE=cmake/arm-none-eabi.cmake
cmake --build build --target sleep_rtc sleep_lptim sleep_tickless -j
probe-rs run --chip STM32H7A3ZITx build/sleep_rtc.elf
python meas/capture.py --seconds 120 --out trace.npz
python meas/cycles.py trace.npz && python meas/battery.py
\end{shellcode}

\begin{note}{When the board sleeps and the debugger does not}
A debugger attached across a sleep state changes what you are measuring and
sometimes prevents the state entirely, because the debug block keeps clocks
running so that it can still see the core. Capture with the debugger detached
and the board reset by a power cycle, and say so in the caption. If the board
will not re-enter debug afterwards, connect under reset. If a build sleeps
immediately and never enumerates, hold the user button at reset and have the
first thing the application does be a check of that button, so there is always a
path back into a board that has decided to be quiet.
\end{note}

\subsection*{Verification and acceptance criteria}

\begin{itemize}
\item The instrument's reading against a known resistor at a known voltage is
recorded, with its deviation, before any board measurement is reported.
\item Charge and energy per cycle are reported as a median with a minimum and a
maximum over at least ten consecutive cycles. A single average current is not
an acceptable result in this chapter or in chapters 10 and 12.
\item The trace with the clock restore and the trace without it are captured
under otherwise identical conditions and both are published. The difference is
stated in microjoules per cycle and in restore milliseconds.
\item Every current figure carries the list of what was enabled when it was
taken: frequency, voltage scale, peripheral clocks, caches, and whether the
debugger was attached.
\item The wake period is checked against the real-time clock over at least an
hour and the drift is stated, which is claimable here because the crystal is
fitted.
\item The three sleep builds share one unmodified application source, and the
analysis passes its synthetic test on the host with no instrument attached.
\item The battery number is printed together with its assumptions in the same
output, so that the two cannot be separated by whoever quotes it.
\end{itemize}

\subsection*{Variants}

\begin{table}[H]\centering\small
\begin{tabular}{p{24mm}p{26mm}p{44mm}p{22mm}p{20mm}}
\toprule
Axis & Variant & What changes & Cost & Built in full in \\
\midrule
Time and safety & The real-time clock & A calendar alarm from the fitted crystal wakes the part. The period is exact and the wall-clock time survives the sleep & The calendar's own setup, and a backup domain question & Here \\
Time and safety & Low-power timer & A counter that continues in the stop state because it runs from the crystal rather than a bus clock. No calendar & Fewer features than a general timer & Here \\
Time and safety & Tickless idle & The millisecond tick is stopped before sleeping and the elapsed time is reconstructed from the counter that kept running & Granularity of one crystal tick & Here \\
Time and safety & General timer or SysTick & Both stop when their clock stops, so neither can wake the part from the stop state at all & Not usable here & Chapter 6 \\
Time and safety & Standby instead of stop & Deeper, and the vendor ships two examples of it for this board. Memory contents are not retained, so the application restarts & A restart is not a wake & Referenced here \\
Time and safety & Watchdog across the sleep & Whether the independent watchdog survives the stop state is item 28 on the confirm list and decides whether a sleeping node can be recovered automatically & A missed wake becomes a reset & Chapter 12 \\
Operating system & RTOS tickless idle & The same reconstruction, owned by the kernel, with the idle task deciding when it is safe & A kernel and its assumptions & Chapter 20 \\
Language & MicroPython & The interpreter exposes sleep states directly, which is the fastest way to see the current fall & No control of the restore & Chapter 1 \\
\bottomrule
\end{tabular}
\caption{Variants for chapter 7. The first three are built here because they are the three that still run when the bus clocks stop, which is the property that defines this chapter.}
\end{table}

The comparison worth making across those rows is not which sleep state is
deepest. It is which one leaves a usable timebase running, because a node that
sleeps perfectly and cannot tell how long it slept is not a node. The real-time
clock keeps a calendar, the low-power timer keeps a count, and the standby state
keeps neither in the general case. Chapter 10 takes the same three and asks a
different question, which is where inside a cycle the charge actually goes.

\subsection*{Pitfalls}

\begin{itemize}
\item Omitting the clock restore after wake and then reporting the energy. The
board works, so nothing points at the cause.
\item Restoring the frequency before the voltage scale and the flash wait
states, in either this chapter or chapter 1. The order is not decorative.
\item Taking a power measurement with the debugger attached and not saying so.
The debug block keeps clocks running and the number is not the node's.
\item Copying a power control register sequence from a project for the
better-known member of this family. This part has a different domain structure
and the bit names differ.
\item Enabling a wake source in its own peripheral but not in the wake-line
register. The interrupt fires and the part stays asleep.
\item Entering the wait instruction with a flag still pending from the previous
cycle, which returns immediately and presents as a board that never sleeps.
\item Writing the low-power timer's autoreload register before enabling the
peripheral, or without waiting for the write to be accepted.
\item Reporting an average current over a duty-cycled trace. It hides whether
the cycles are alike, which is the only thing a battery cares about.
\item Quoting a battery lifetime without the assumptions it rests on.
\end{itemize}

\subsection*{Best practices applied}

\begin{itemize}
\item The finding comes first. The chapter states that this part is a poor
low-power node before it states anything else, and every claim after that is
framed as method rather than as a result to be proud of.
\item Every measurement names its instrument, and every current figure carries
the configuration it was taken under.
\item A dependency's licence is stated before any code is written around it.
The instrument's interface is GPL-2.0, so it is installed and pinned and never
vendored.
\item Measurement and arithmetic are kept apart. The energy per cycle is
measured; the lifetime is computed, and the assumptions travel with it.
\item No figure in microamps appears anywhere until an instrument has produced
it, including in the budget table.
\item The spread is reported with the median, because a duty-cycled system that
is usually consistent and occasionally not is a different system from one that
is always consistent.
\end{itemize}

\subsection*{Stretch goals}

\begin{itemize}
\item Find the crossover. Sweep the sense-phase duration and the wake period and
find the point at which racing to sleep on a fast part stops paying, then state
it as a rule with its conditions. Nobody has published that curve for this
family.
\item Repeat the whole measurement on an ultra-low-power board if one ever
reaches this bench, using the same scripts, and publish the two together. The
method transfers even though the numbers do not.
\item Measure the restore time as a function of the phase-locked loop
configuration and find the cheapest configuration that still meets the
application's throughput requirement.
\item Add a fourth build that wakes on the user button rather than on time, and
show what an event-driven node costs when the events are rare.
\item Write the whole thing up as the measurement method it is, since no
peer-reviewed method for this instrument exists, and offer it for review rather
than only publishing numbers.
\end{itemize}

\subsection*{Roadmap and next steps}

Chapter 8 gives the sleeping node something to do between wakes, which is the
state machine this chapter's duty cycle is a skeleton of, and chapter 10 splits
the cycle measured here into its phases and asks where the charge goes.

For going deeper, the community embedded roadmap is the map worth following,
and its point that projects settle questions reading leaves open applies with
force here, because the published material on this specific question is thin.
The free course from a different silicon vendor on its own development kit is
the best free treatment of duty-cycled node design in the pool, and its
concepts transfer directly even though its silicon does not. For the standards a
reader will eventually meet on a product that claims a battery life, the
consumer device security baseline and the industrial component security series
are listed in Appendix H. Nothing in the pool teaches the measurement rigour
this chapter needs, which is why the chapter writes it.

\subsection*{Portfolio evidence}

\begin{itemize}
\item A repository with three sleep builds, one application source, and a
README that states the finding about this part in its first paragraph rather
than burying it.
\item Two published traces under otherwise identical conditions, with and
without the clock restore, and the difference stated in microjoules per cycle.
\item A per-cycle table over at least ten cycles with the median and both
extremes, and the configuration list attached to it.
\item A battery number printed together with the assumptions it rests on, in the
same output, so it cannot be quoted alone.
\end{itemize}

\subsection*{Sources}

Normative references:
\begin{itemize}
\item Reference manual RM0455, for the low-power modes, what each retains, the
wake sources, the domain structure and the power control register bit names.
Not RM0433, whose domain structure differs.
\item The STM32H7A3xI datasheet, for the voltage scales, the clock limits at
each scale and the supply configuration.
\item The board user manual for MB1363, for the measurement jumper, what it
isolates, and whether the probe's current is excluded from it.
\item The nRF-PPK2 user guide, for the two modes, the maximum current, the
digital input count and threshold, the sample rate and the export format.
\end{itemize}

Reusable implementations:
\begin{itemize}
\item The vendor's three power examples for this exact board, under
\texttt{Projects/}\allowbreak\texttt{NUCLEO-H7A3ZI-Q/}\allowbreak\texttt{Examples/PWR}: standby, standby with a clock
wake, and the deeper stop mode with a clock wake. Basic examples are
BSD-3-Clause; check each file's header.
\item The Python interface to the power instrument, GPL-2.0. Install from the
package index and do not vendor it.\newline
\url{https://github.com/IRNAS/ppk2-api-python}
\item Golioth's article on measuring current consumption with this
instrument, which is the method this chapter starts from.\newline
\url{https://blog.golioth.io/measuring-current-consumption-with-power-profiler-kit-ii/}
\item The MicroPython board definition for this exact board, MIT, which is the
second independent confirmation that the low-speed crystal is fitted.\newline
\url{https://github.com/micropython/micropython}
\end{itemize}
